% GENERATED FILE. Edit canonical lessons, then run npm run generate:books.
% canonical-source-hash: fnv1a64:664789ddc7990979
% canonical-lessons: JA-W01-i, JA-W01-ha, JA-W01-hai-read, JA-W01-e, JA-W01-ko, JA-W01-n, JA-W01-ni, JA-W01-chi, JA-W01-wa, JA-W01-konnichiwa-read

\chapter{Eight Signs, and Three Words You Can Read}
\label{ch:ja-hiragana-eight}
\hlchaptermodality{2}

\begin{chapteropening}
\textbf{By the end of this chapter:} \emph{I can write eight hiragana signs from memory and read the daytime greeting and both yes-or-no answers by sounding out their signs in order, including the sign that is pronounced wa because it is a topic marker.}

The daytime greeting assembled from five signs the reader can each write from memory, introducing nothing new — including the one read wa because it is a topic marker rather than the sign for that sound.
\end{chapteropening}

\section[i]{\ja{い} — two strokes, one beat}

\label{lesson:JA-W01-i}



\begin{quote}
You already know the rule that makes this writing system easy: \textbf{one sign, one beat.} Not one letter one sound, and not one sign one word — one sign, one \textbf{mora}, the even beat Japanese counts time in.

What you have not had yet is a single sign taught on its own. Here is the first.
\end{quote}

\subsection*{Script — the shape}
\begin{quote}
\ja{い}
\end{quote}

Two strokes, both falling, the left one long and the right one short:

1. the \textbf{left} stroke — start high, come down, and hook gently to the right at the bottom 2. the \textbf{right} stroke — a short downward tick, higher up and clear of the first

It is read \textbf{i}, the vowel of English \emph{machine} — short, tight, never the \emph{eye} of English \emph{bite}.

\textbf{Where the shape came from, and why it does not help you.} Every hiragana sign began as a Chinese character written fast and cursively until it wore smooth. But those characters were chosen \textbf{for their sound}, not their meaning, so unlike a Chinese component there is nothing inside \ja{い} to decode. The shape is a signature you learn, not a structure you read. This book will tell you when a shape carries information and when it does not, and this is the second kind.

\subsection*{Your turn}
\begin{itemize}
\raggedright
  \item \emph{Write it:} \ja{い} — long stroke first, then the short tick
  \item \emph{Say it:} \textbf{i}, one clean beat
\end{itemize}

\begin{tcolorbox}[breakable,colback=teal!4,colframe=teal!35!black,title={Before you move on}]
How many beats does one hiragana sign take? (\textbf{One.})

Source: \href{https://www.unicode.org/charts/PDF/U3040.pdf}{Unicode Hiragana chart}; stroke order per \emph{Hitsujun Shido no Tebiki} (Japanese Ministry of Education, 1958).
\end{tcolorbox}
\section[ha]{\ja{は} — three strokes}

\label{lesson:JA-W01-ha}



\begin{quote}
Write the two-stroke sign you met a moment ago. Long stroke, short tick.
\end{quote}

\subsection*{Script — the shape}
\begin{quote}
\ja{は}
\end{quote}

Three strokes. Look at it as a tall left post with a loop tied on its right.

1. the \textbf{vertical}, down the left side, with the smallest flick left at the foot 2. a \textbf{short horizontal} crossing that vertical near the top, on its right 3. the \textbf{loop} — down from the right of that crossbar, round to the left, and back up to close

Read \textbf{ha}, breathed rather than hard: the \emph{h} of English \emph{hat}, then the open \emph{a} of \emph{father}.

\textbf{It has a second reading, and you have already used it.} Keep that in your pocket for now — this sign is doing two jobs in Japanese, and the second one is what makes the greeting look misspelled to a beginner. It gets its own lesson, after you have met the sign it \emph{sounds} like.

\subsection*{Your turn}
\begin{itemize}
\raggedright
  \item \emph{Write it:} \ja{は} — vertical, crossbar, loop
  \item \emph{Say it:} \textbf{ha}, one beat, breath not force
\end{itemize}

\begin{tcolorbox}[breakable,colback=teal!4,colframe=teal!35!black,title={Before you move on}]
Which stroke of \ja{は} is written first? (The \textbf{vertical} on the left.)

Source: \href{https://www.unicode.org/charts/PDF/U3040.pdf}{Unicode Hiragana chart}; stroke order per \emph{Hitsujun Shido no Tebiki} (Japanese Ministry of Education, 1958).
\end{tcolorbox}
\section[hai]{Two signs, side by side — and a word appears}

\label{lesson:JA-W01-hai-read}



\begin{quote}
Both signs are yours. Write them, in this order, and say each beat as you finish it:

\begin{quote}
\ja{は}   \ja{い}
\end{quote}
\end{quote}

\subsection*{Script — no new sign at all}
Now write them touching:

\begin{quote}
\textbf{\ja{は}} + \textbf{\ja{い}} $\to$ \textbf{\ja{はい}}
\end{quote}

\textbf{ha} then \textbf{i}. Two beats, even, neither one stressed. That is the word for \textbf{yes}, and it is the one you have been saying since your first minute in this language.

\textbf{Nothing new was introduced in this lesson.} No stroke, no sound, no word. What changed is where the word now lives: you are no longer recalling a shape you were shown, you are \textbf{reading} it — assembling it out of two pieces you can write from memory and sound out in order.

That is the whole method for this writing system, and it does not get harder. It gets longer. A word of five signs is five signs you already own, in a row.

\subsection*{Your turn}
\begin{itemize}
\raggedright
  \item \emph{Read it:} \textbf{\ja{はい}} — sound each sign in turn, do not recall the word whole
  \item \emph{Write it:} \ja{はい} — five strokes in total, three then two
\end{itemize}

\begin{tcolorbox}[breakable,colback=teal!4,colframe=teal!35!black,title={Before you move on}]
How many new strokes did you have to learn in this lesson? (\textbf{None.})

Source: \href{https://www.unicode.org/charts/PDF/U3040.pdf}{Unicode Hiragana chart}.
\end{tcolorbox}
\section[e]{\ja{え} — two strokes, and a second word for free}

\label{lesson:JA-W01-e}



\begin{quote}
Read the word for \emph{yes} off the page, sign by sign, out loud.
\end{quote}

\subsection*{Script — the shape}
\begin{quote}
\ja{え}
\end{quote}

Two strokes, and the second one does most of the work:

1. a \textbf{short tick} at the top, on its own, clear of everything below 2. the \textbf{body} — across to the right, back down and left in a flat zigzag, then out to the lower right in one continuous movement

Read \textbf{e}, the vowel of English \emph{bed} — short, and never the \emph{ay} of \emph{day}.

\begin{grammarlens}[title={What you've built}]
Three signs are now yours, and they spell more than one word. Put two of them together with the one you just learned:

\begin{quote}
\textbf{\ja{い}} + \textbf{\ja{い}} + \textbf{\ja{え}} $\to$ \textbf{\ja{いいえ}}
\end{quote}

\textbf{i-i-e.} Three beats — and note that the first sign is written \textbf{twice}, because Japanese counts two full beats there. That is the word for \textbf{no}.

You have now read both answers to a yes-or-no question, from the signs up. Neither of them was memorised as a picture.
\end{grammarlens}

\subsection*{Your turn}
\begin{itemize}
\raggedright
  \item \emph{Write it:} \ja{え} — tick first, then the body in one movement
  \item \emph{Read it:} \ja{いいえ} — three beats, and do not shorten the doubled sign
\end{itemize}

\begin{tcolorbox}[breakable,colback=teal!4,colframe=teal!35!black,title={Before you move on}]
Why is the first sign of that word written twice? (Because it is \textbf{two beats}, not one held longer.)

Source: \href{https://www.unicode.org/charts/PDF/U3040.pdf}{Unicode Hiragana chart}; stroke order per \emph{Hitsujun Shido no Tebiki} (Japanese Ministry of Education, 1958).
\end{tcolorbox}
\section[ko]{\ja{こ} — two strokes, and the greeting begins}

\label{lesson:JA-W01-ko}



\begin{quote}
The two-stroke sign read \emph{e} — tick, then body. Write it once.
\end{quote}

\subsection*{Script — the shape}
\begin{quote}
\ja{こ}
\end{quote}

The simplest sign you have met. Two short horizontal strokes, stacked, the lower one longer and slightly curved:

1. the \textbf{upper} stroke, short, curving down a little at its right end 2. the \textbf{lower} stroke, wider, sweeping right and lifting slightly

Read \textbf{ko}: a clean \emph{k}, then the \emph{o} of English \emph{more} shortened.

Nothing joins them. The gap between the two strokes is part of the shape — close it and you have a different sign.

\subsection*{Your turn}
\begin{itemize}
\raggedright
  \item \emph{Write it:} \ja{こ} — upper stroke, then the longer lower one, and leave the gap
  \item \emph{Say it:} \textbf{ko}, one beat
\end{itemize}

\begin{tcolorbox}[breakable,colback=teal!4,colframe=teal!35!black,title={Before you move on}]
Which of the two strokes is longer? (The \textbf{lower} one.)

Source: \href{https://www.unicode.org/charts/PDF/U3040.pdf}{Unicode Hiragana chart}; stroke order per \emph{Hitsujun Shido no Tebiki} (Japanese Ministry of Education, 1958).
\end{tcolorbox}
\section[n]{\ja{ん} — one stroke, one beat, no vowel}

\label{lesson:JA-W01-n}



\begin{quote}
Two stacked strokes with a gap. Write the sign read \emph{ko}.
\end{quote}

\subsection*{Script — the shape, and the exception}
\begin{quote}
\ja{ん}
\end{quote}

\textbf{One stroke.} Start with a small tick down-left, then sweep down, round to the right and up in a single unbroken movement, finishing with a rise.

Now the part that matters more than the shape. Every sign you have met so far is a consonant plus a vowel, or a vowel alone. \textbf{This one is a consonant with no vowel at all} — and it still takes a \textbf{full beat}, exactly as long as any other sign.

That is why Japanese counts morae rather than syllables, and it is the clearest evidence you will get. An English speaker hears the \emph{n} at the end of a word as a tail on the vowel before it. Japanese does not. It is its own beat, held its own length, and shortening it is one of the most audible foreign-accent tells there is.

\subsection*{Your turn}
\begin{itemize}
\raggedright
  \item \emph{Write it:} \ja{ん} — one continuous stroke, no lifts
  \item \emph{Say it:} hold it for a full beat, as long as any other sign takes
\end{itemize}

\begin{tcolorbox}[breakable,colback=teal!4,colframe=teal!35!black,title={Before you move on}]
If this sign has no vowel, how long is it held? (\textbf{A full beat} — the same as every other sign.)

Source: \href{https://www.unicode.org/charts/PDF/U3040.pdf}{Unicode Hiragana chart}; stroke order per \emph{Hitsujun Shido no Tebiki} (Japanese Ministry of Education, 1958).
\end{tcolorbox}
\section[ni]{\ja{に} — three strokes, and a shape you half know}

\label{lesson:JA-W01-ni}



\begin{quote}
Write the three-stroke sign read \emph{ha}. Vertical, crossbar, loop.
\end{quote}

\subsection*{Script — the shape}
\begin{quote}
\ja{に}
\end{quote}

Three strokes, and \textbf{the first one you have already practised}:

1. the \textbf{vertical} down the left, with the same small flick left at the foot as the sign read \emph{ha} 2. a \textbf{short horizontal} to its right, upper area 3. a \textbf{longer horizontal} below it, sweeping right

Read \textbf{ni}: the \emph{n} of \emph{knee}, then the tight \emph{i} you learned first.

\textbf{This is the first time two signs share a visible part}, and it is worth naming, because it is nearly the last. Hiragana is not built from recombining components the way Chinese characters are — these two happen to have worn down to similar left-hand posts, not because they mean or sound anything alike. Treat the resemblance as a \textbf{handwriting} convenience, not as a clue. Reading information into it is the mistake this book keeps warning you off.

\subsection*{Your turn}
\begin{itemize}
\raggedright
  \item \emph{Write it:} \ja{に} — vertical, short bar, long bar
  \item \emph{Notice:} write it beside \ja{は} and compare only the left post
\end{itemize}

\begin{tcolorbox}[breakable,colback=teal!4,colframe=teal!35!black,title={Before you move on}]
Which stroke of this sign comes first? (The \textbf{vertical}.)

Source: \href{https://www.unicode.org/charts/PDF/U3040.pdf}{Unicode Hiragana chart}; stroke order per \emph{Hitsujun Shido no Tebiki} (Japanese Ministry of Education, 1958).
\end{tcolorbox}
\section[chi]{\ja{ち} — two strokes, four roman letters, one beat}

\label{lesson:JA-W01-chi}



\begin{quote}
Three strokes: vertical, short bar, long bar. Write the sign read \emph{ni}.
\end{quote}

\subsection*{Script — the shape}
\begin{quote}
\ja{ち}
\end{quote}

Two strokes:

1. a \textbf{short horizontal} across the top 2. the \textbf{body} — down from the middle of that bar, curving left, then round into a full belly to the right and up

Read \textbf{chi}: the \emph{ch} of English \emph{cheese}, then the tight \emph{i}.

\textbf{Count what just happened.} One sign. One beat. \textbf{Four roman letters} to describe it. The romanization is a pronunciation aid written for people who do not read the script yet — it is not a smaller version of the writing system, and the number of letters in it tells you nothing about the number of beats in the word.

This is the moment the roman spelling starts costing you more than it pays, and it is why the signs are worth the effort.

\subsection*{Your turn}
\begin{itemize}
\raggedright
  \item \emph{Write it:} \ja{ち} — top bar, then the body in one movement
  \item \emph{Say it:} \textbf{chi}, one beat — as long as any other sign, no longer
\end{itemize}

\begin{tcolorbox}[breakable,colback=teal!4,colframe=teal!35!black,title={Before you move on}]
Does the number of roman letters tell you how long a sign is held? (\textbf{No.})

Source: \href{https://www.unicode.org/charts/PDF/U3040.pdf}{Unicode Hiragana chart}; stroke order per \emph{Hitsujun Shido no Tebiki} (Japanese Ministry of Education, 1958).
\end{tcolorbox}
\section[wa]{\ja{わ} — the sign for \emph{wa}, and the trap beside it}

\label{lesson:JA-W01-wa}



\begin{quote}
Write the three-stroke sign read \emph{ha}: vertical, crossbar, loop.
\end{quote}

\subsection*{Script — the shape}
\begin{quote}
\ja{わ}
\end{quote}

Two strokes:

1. the \textbf{vertical} on the left, with the small flick at the foot 2. the \textbf{body} — a short bar right from the top, then down, round to the left in a loop, and out to the lower right with a curl

Read \textbf{wa}: like English \emph{wa} in \emph{wander}, one beat.

\begin{grammarlens}[title={What you've built — and the trap}]
Here is the thing this sign exists in this lesson to protect you from.

The daytime greeting \textbf{sounds} like it ends in \emph{wa}. \textbf{It is not written with this sign.} It ends with the sign read \emph{ha} — the one you learned early and were told was holding a second job.

That second job is grammatical: \ja{は} is also Japanese's topic marker, and \textbf{as a topic marker it is pronounced \emph{wa}}. The greeting is a frozen fragment of a longer sentence, and it kept the marker's spelling along with the marker's sound.

So the rule you need is not about sound at all:

\begin{itemize}
\raggedright
  \item \textbf{\ja{わ}} is the sign for the mora \emph{wa} inside ordinary words.
  \item \textbf{\ja{は}} read as \emph{wa} means you are looking at the topic marker, or at something that used to contain one.
\end{itemize}

Getting this backwards is the single most common beginner spelling error in Japanese, and it is the reason this book taught you the sign for \emph{wa} \textbf{after} you already knew the greeting — so the shape could arrive with its warning attached rather than be quietly assumed.
\end{grammarlens}

\subsection*{Your turn}
\begin{itemize}
\raggedright
  \item \emph{Write it:} \ja{わ}, then \ja{は} beside it — two signs, one shared sound, different jobs
  \item \emph{Decide:} hearing \emph{wa} at the end of the daytime greeting, which sign do you write?
\end{itemize}

\begin{tcolorbox}[breakable,colback=teal!4,colframe=teal!35!black,title={Before you move on}]
Which of the two is a fact about \textbf{grammar} rather than about sound? (The one where \ja{は} is read \emph{wa}.)

Source: \href{https://www.unicode.org/charts/PDF/U3040.pdf}{Unicode Hiragana chart}; stroke order per \emph{Hitsujun Shido no Tebiki} (Japanese Ministry of Education, 1958).
\end{tcolorbox}
\section[konnichiwa]{Five signs in a row — and the greeting is readable}

\label{lesson:JA-W01-konnichiwa-read}



\begin{quote}
Four signs, one after another, from memory. Say each beat as you finish writing it:

\begin{quote}
\ja{こ}   \ja{ん}   \ja{に}   \ja{ち}
\end{quote}
\end{quote}

\subsection*{Script — the assembly}
Add the fifth — the one that is read \emph{wa} because it is a topic marker — and the greeting is on the page:

\begin{quote}
\textbf{\ja{こ}} + \textbf{\ja{ん}} + \textbf{\ja{に}} + \textbf{\ja{ち}} + \textbf{\ja{は}} $\to$ \textbf{\ja{こんにちは}}
\end{quote}

\textbf{ko-n-ni-chi-wa.} Five signs, \textbf{five beats}, all even. Say it against a steady tap and you will hear that the second beat, the one with no vowel, takes exactly as long as the four around it.

\textbf{Nothing new was introduced in this lesson.} You have not learned a stroke, a sound or a word here. You have read a word you already knew how to say — from five pieces you can each write from memory, in order, including the one that sounds like something other than its name.

\begin{grammarlens}[title={What you've built}]
Eight signs, taught one at a time, and with them you can now read three whole words off the page: the greeting, and both answers to a yes-or-no question.

Hiragana has forty-six basic signs. You hold \textbf{eight}. The rest arrive the same way, one per lesson, and the words they unlock arrive free — because a word in this script is never more than the signs you already own, in a row.
\end{grammarlens}

\subsection*{Your turn}
\begin{itemize}
\raggedright
  \item \emph{Read it:} \textbf{\ja{こんにちは}} — sound each sign in turn; do not recall it whole
  \item \emph{Write it:} the greeting, then say it once at speaking pace
\end{itemize}

\begin{tcolorbox}[breakable,colback=teal!4,colframe=teal!35!black,title={Before you move on}]
How many of the greeting's five signs were new in this lesson? (\textbf{None.})

Source: \href{https://www.unicode.org/charts/PDF/U3040.pdf}{Unicode Hiragana chart}.
\end{tcolorbox}
